%% 5-1 平均の速さと瞬間の速さ（第11回：ページ制限なし。偶数ページ）
\bkchapter{5}{1}{平均の速さと瞬間の速さ}{\bklevel{2}}{平均の速さと瞬間の速さ}
\bkmichishirube{2-3,3-4,3-5}{5-2,5-3}

\begin{goal}
\begin{itemize}[label={\color{bkTheme}●}, leftmargin=1.6\zw, itemsep=0pt, topsep=1pt]
\item \bklevel{2}\enspace 平均の速さを、表やグラフから求められる。
\item \bklevel{2}\enspace 平均の速さが、グラフの「2点を結ぶ直線の傾き」になることを説明できる。\bkidx{かたむき}{傾き}
\item \bklevel{2}\enspace 時間を短くして、瞬間の速さを見積もれる。
\end{itemize}
\end{goal}

\begin{hatena}
\qty{100}{km}はなれた町まで、車で2時間。$100\div2=50$ で時速\qty{50}{km}です。でも、ずっと時速\qty{50}{km}で走っていたのでしょうか？\par
\begin{bkfigure}
\begin{tikzpicture}
% 第7回：区間をはっきり分けた折れ線にする（0〜0.5h 時速60、0.5〜0.7h 止まる、0.7〜1h 時速50、
%   1〜1.25h じゅうたい時速20、1.25〜1.5h 時速80、1.5〜2h 時速60。合計100 km）。
%   0.5時間ごとの点は例題1の表と同じ。注記は帯の上と右下の凡例に置き、引き出し線は使わない
\begin{axis}[bkplot, width=100mm, height=50mm, xmin=0, xmax=2.2, ymin=0, ymax=112,
  xtick={0.5,1,1.5,2}, ytick={25,50,75,100}, xlabel={時間（h）}, ylabel={進んだ距離（km）},
  ylabel style={anchor=south east, xshift=3mm}]
  \fill[bkRed!12] (axis cs:0.5,0) rectangle (axis cs:0.7,104);
  \fill[bkOrange!16] (axis cs:1.0,0) rectangle (axis cs:1.25,104);
  \node[font=\sffamily\scriptsize, text=bkRed!80!black, anchor=south, align=center, inner sep=1pt] at (axis cs:0.6,104) {止まった\\（休けい）};
  \node[font=\sffamily\scriptsize, text=bkOrange!80!black, anchor=south, align=center, inner sep=1pt] at (axis cs:1.125,104) {じゅうたい\\（のろのろ）};
  \addplot[bkthin, dash pattern=on 4pt off 2pt, domain=0:2] {50*x};
  \addplot[bkcurve, line join=round] coordinates {(0,0) (0.5,30) (0.7,30) (1.0,45) (1.25,50) (1.5,70) (2,100)};
  \addplot[only marks, mark=*, mark size=1.5pt, bkInk] coordinates {(0.5,30) (1,45) (1.5,70) (2,100)};
  \node[on layer=axis foreground, anchor=south east, fill=white, draw=bkLine, line width=0.4pt, rounded corners=1mm, inner sep=1.2mm,
    font=\sffamily\scriptsize, text=bkInk, align=left] at (axis cs:2.17,4)
    {\tikz[baseline=-0.6ex]\draw[bkBlue, line width=1.7pt] (0,0)--(5mm,0);\enspace じっさいの車\\[1pt]
     \tikz[baseline=-0.6ex]\draw[bkGray, line width=0.9pt, dash pattern=on 4pt off 2pt] (0,0)--(5mm,0);\enspace ずっと時速\qty{50}{km}};
\end{axis}
\end{tikzpicture}
\end{bkfigure}
水平な所は、信号や休けいで止まっていた所（メーターは0）。傾きがゆるい所は、じゅうたいの所（メーターは20くらい）。「時速\qty{50}{km}」と、スピードメーターが指す速さは、何がちがうのでしょう。
\end{hatena}

\clearpage
%% ---------------- やってみよう・思い出そう ----------------
\begin{yattemiyou}[坂をころがるボールの速さを調べよう]
\textbf{用意するもの}\par
長い板（または本を重ねてつくった坂）、ボール、ものさし、スマートフォン（動画をとる）
\begin{enumerate}[label={\bkcircled[bkGreen]{\arabic*}}, leftmargin=2.4\zw, itemsep=1pt, topsep=3pt]
\item 板を少しだけ傾け、ものさしを横に置く。ボールを静かに放して、ころがるようすを動画にとる。
\item 動画を0.5秒ずつ止めて、ボールの位置を読み、表に書く。
\item となりどうしの時刻のあいだで、「進んだ距離÷かかった時間」を計算しよう。
\end{enumerate}
\begin{bkcenter}
\begin{bktable}{colspec={Q[l,m,wd=34mm] *{5}{Q[c,m,wd=13mm]}}, rowsep=1.5pt}
時刻 $t$（s） & 0 & 0.5 & 1.0 & 1.5 & 2.0 \\
位置 $x$（cm） & 0 & 5 & 20 & 45 & 80 \\
\end{bktable}
\bktabcaption{ある班の結果。坂の上のスタート地点を0にした。}
\end{bkcenter}
\end{yattemiyou}

\noindent ③の計算をしてみましょう。0.5秒ごとに進んだ距離は、「あとの位置${}-{}$まえの位置」です。それを、かかった時間0.5秒でわります。

\begin{bkcenter}
\begin{bktable}{colspec={Q[l,m,wd=27mm] *{4}{Q[c,m,wd=23.5mm]}}, rowsep=1.5pt, colsep=2pt, cells={font=\small}}
あいだ（s） & 0〜0.5 & 0.5〜1.0 & 1.0〜1.5 & 1.5〜2.0 \\
進んだ距離（cm） & $5-0=5$ & $20-5=15$ & $45-20=25$ & $80-45=35$ \\
距離÷時間（cm/s） & $5\div0.5=10$ & $15\div0.5=30$ & $25\div0.5=50$ & $35\div0.5=70$ \\
\end{bktable}
\bktabcaption{0.5秒ごとの「距離÷時間」}
\end{bkcenter}

なぜ0.5でわるのでしょう。0.5秒で\qty{5}{cm}進んだボールが、同じ速さで1秒（0.5秒の2倍）走れば、2倍の\qty{10}{cm}進みます。だから「1秒あたり\qty{10}{cm}」で、\qty{10}{cm/s}です。0.5でわることは、2倍することと同じです（$5\div0.5=10$）。

「距離÷時間」は、10、30、50、70と増えています。ボールは、だんだん速くなっています。では、ちょうど1.0秒の「瞬間」の速さは、いくつなのでしょう。この章では、それを考えます。

\begin{omoidasou}
\bkrecall{2-3}{速さ${}={}$進んだ道のり÷かかった時間。変わった量を \mbox{$\Delta$（デルタ）}で書く。$\Delta x$ は「$x$ の変化量」。}
\bkrecall{3-4}{横軸を時間 $t$、たて軸を位置 $x$ にしたグラフを、$x$-$t$ グラフという。}
\bkrecall{3-5}{グラフの傾き $=\dfrac{\text{たての変化量}}{\text{横の変化量}}$。直線が右上がりなら正、右下がりなら負。}
\bkrecall{2-10}{\qty{1000}{m/km}（1 km あたり 1000 m）、\qty{3600}{s/h}（1時間あたり3600秒）。\\換算は、単位つきの数をかけたりわったりして、同じ単位を約分して消す。}
\end{omoidasou}

%% ---------------- 数学の道具 ----------------
\bksection{doguu}{数学の道具}

\bksubsection{1}{平均の速さは $\dfrac{\Delta x}{\Delta t}$}

\label{bk:5-1:heikin}\bkhatenaref の車で考えます。出発して0.5時間（30分）のとき、車は出発点から\qty{30}{km}の所にいました。1時間のときは\qty{45}{km}の所にいました。このあいだの速さを求めましょう。

\begin{bkfigure}
\begin{tikzpicture}[x=1mm, y=1mm]
  % 道路を数直線にする（1 km＝1.1 mm）。0.5時間の車と1時間の車。Δx は 30→45 km
  \draw[bkGray, line width=0.8pt] (0,0) -- (118,0);
  \foreach \k in {0,10,...,100} {\draw[bkGray, line width=0.6pt] (\k*1.1,0) -- (\k*1.1,-1.8);
    \node[anchor=north, font=\sffamily\scriptsize] at (\k*1.1,-2.4) {\k};}
  \node[anchor=west, font=\sffamily\scriptsize] at (113,-4.2) {（km）};
  \bkcar{(30*1.1-11.3,0.3)}{0.42}{bkBlue!70}
  \bkcar{(45*1.1-11.3,0.3)}{0.42}{bkOrange}
  \node[anchor=south, font=\sffamily\scriptsize, text=bkBlue] at (30*1.1-4,8.5) {0.5時間のとき};
  \node[anchor=south, font=\sffamily\scriptsize, text=bkOrange!80!black] at (45*1.1+3,8.5) {1時間のとき};
  \draw[bkInk, line width=0.6pt, {Bar[width=4pt]}-{Bar[width=4pt]}] (30*1.1,-9.5) -- (45*1.1,-9.5);
  \node[anchor=north, font=\sffamily\footnotesize] at (37.5*1.1,-10.8) {$\Delta x=45-30=\qty{15}{km}$};
  \node[anchor=north west, font=\sffamily\footnotesize, align=left] at (66,-8.5) {かかった時間\\$\Delta t=1-0.5=\qty{0.5}{h}$};
\end{tikzpicture}
\bkfigcaption{道路を数直線にして、0.5時間と1時間の車の位置をかいた。}
\end{bkfigure}

進んだ距離もかかった時間も、「あと${}-{}$まえ」の変化量です（2-3）。位置 $x$ の変化量を $\Delta x$、時刻 $t$ の変化量を $\Delta t$ と書きます。\bkidx{へんかりょう}{変化量}
\[ \Delta x=45-30=15\quad\text{（km）},\qquad \Delta t=1-0.5=0.5\quad\text{（h）} \]
速さは、距離を時間でわって
\[ \frac{\Delta x}{\Delta t}=\frac{15}{0.5}=30\quad\text{（km/h）} \]
\textbf{時速\qty{30}{km}}です。0.5でわるのは、やってみようと同じで、「0.5時間で\qty{15}{km}なら、1時間ではその2倍」だからです。

この時速\qty{30}{km}は、「このあいだ、ずっと時速\qty{30}{km}で走ったとすると、ちょうど同じ\qty{15}{km}進む」という速さです。じっさいには、このあいだに車は止まったり（休けい）、時速\qty{50}{km}で走ったりしています。それを\textbf{ならした}値なので、\textbf{平均の速さ}といいます。\bkidx{へいきんのはやさ}{平均の速さ}

いつでも使えるように、文字で書きます。時刻 $t_1$ に位置 $x_1$ にいて、時刻 $t_2$ に位置 $x_2$ にいたとします。$t_1$ の右下の小さな1は「1つめ」という印で、「ティー・いち」と読みます。かけ算ではありません。
\[ \Delta t=t_2-t_1,\qquad \Delta x=x_2-x_1,\qquad \text{平均の速さ}=\frac{\text{進んだ距離}}{\text{かかった時間}}=\frac{\Delta x}{\Delta t} \]

同じ車でも、0時間から2時間までをとると、$\Delta x=100-0=100$、$\Delta t=2-0=2$ なので、平均の速さは $100\div2=\qty{50}{km/h}$ です。0.5時間から1時間までの\qty{30}{km/h}とは、ちがう値です。平均の速さは、\textbf{どのあいだをとるか}で変わります。

\bksubsection{2}{平均の速さは、2点を結ぶ直線の傾き}

\label{bk:5-1:katamuki}やってみようのボールの $x$-$t$ グラフに、時刻1.0秒の点Pと、時刻2.0秒の点Qをとり、PとQを結ぶ直線を引きます。

\begin{bkfigure}
\begin{tikzpicture}
\begin{axis}[bkplot, width=78mm, height=48mm, xmin=0, xmax=2.3, ymin=0, ymax=95,
  xtick={0,0.5,1,1.5,2}, ytick={0,20,40,60,80}, xlabel={$t$（s）}, ylabel={$x$（cm）}]
  \addplot[bkcurve, domain=0:2.1] {20*x^2};
  \addplot[bkcurve2, domain=0.7:2.2, line width=1.2pt] {20+60*(x-1)};
  \addplot[only marks, mark=*, mark size=2pt, bkInk] coordinates {(1,20) (2,80)};
  \draw[bkGray, line width=0.7pt, dash pattern=on 2pt off 1.5pt] (axis cs:1,20) -- (axis cs:2,20) -- (axis cs:2,80);
  \node[anchor=north, font=\sffamily\footnotesize] at (axis cs:1.5,20) {$\Delta t=\qty{1.0}{s}$};
  \node[anchor=west, font=\sffamily\footnotesize] at (axis cs:2,50) {$\Delta x=\qty{60}{cm}$};
  \node[anchor=south east, font=\sffamily\small] at (axis cs:1,20) {P};
  \node[anchor=south east, font=\sffamily\small] at (axis cs:2,80) {Q};
\end{axis}
\end{tikzpicture}
\bkfigcaption{やってみようのボール。PとQを結ぶ直線（オレンジ）の傾きは $\dfrac{60}{1.0}=60$。}
\end{bkfigure}

直線の傾きは、「たての変化量÷横の変化量」でした（3-5）。$x$-$t$ グラフでは、たては位置 $x$、横は時刻 $t$ なので、
\[ \text{傾き}=\frac{\text{たての変化量}}{\text{横の変化量}}=\frac{x_2-x_1}{t_2-t_1}=\frac{\Delta x}{\Delta t} \]
これは、平均の速さとまったく同じ式です。PからQまでは、たてに $80-20=\qty{60}{cm}$、横に $2.0-1.0=\qty{1.0}{s}$ なので、傾きは $60\div1.0=60$。1.0秒から2.0秒までの平均の速さは\qty{60}{cm/s}です。

もう1つの見方もあります。ずっと同じ速さで走る物の $x$-$t$ グラフは、まっすぐな直線になります（1秒ごとに同じ距離ずつ進むから）。直線PQは、「Pの時刻にPの位置を出て、ずっと同じ速さで走り、Qの時刻にQの位置に着く物」のグラフです。その速さが平均の速さです。\bkhatenaref のグラフの点線（ずっと時速\qty{50}{km}）も、出発の点と着いた点を結ぶ直線になっていました。

このように、$x$-$t$ グラフの上の\textbf{2点を結ぶ直線の傾き}が、そのあいだの平均の速さになります。

\bksubsection{3}{時間を短くすると、瞬間の速さに近づく}

\label{bk:5-1:shunkan}ボールの1.0秒の瞬間の速さを知りたいとします。1.0秒から2.0秒までの平均（\qty{60}{cm/s}）では、あいだが長すぎます。そのあいだにも、ボールはどんどん速くなっているからです。

それなら、あいだを0にすればよいのでしょうか。1.0秒から1.0秒までにすると、進んだ距離は0、かかった時間も0なので、$0\div0$ になって計算できません（1-14）。そこで、あいだを\textbf{0にはしないで、どんどん短くしていく}ことにします。

短いあいだの位置を知るには、表にない時刻の位置が必要です。表をよく見ると、時刻が0.5秒から1.0秒へ2倍になると、位置は5から20へ4倍（$2\times2$）、1.5秒（0.5秒の3倍）では45で9倍（$3\times3$）です。位置は時刻の2乗に比例しています（3-9）。比例定数は、1.0秒のときの $20\div1.0^2=20$ なので、
\[ x=20t^2\quad\text{（cm）} \]
という式で表せます。表の数と合うか、確かめましょう。

\begin{bkcenter}
\begin{bktable}{colspec={Q[l,m,wd=34mm] *{5}{Q[c,m,wd=14mm]}}, rowsep=1.5pt}
時刻 $t$（s） & 0 & 0.5 & 1.0 & 1.5 & 2.0 \\
$t^2$ & 0 & 0.25 & 1 & 2.25 & 4 \\
位置 $x=20t^2$（cm） & 0 & 5 & 20 & 45 & 80 \\
\end{bktable}
\bktabcaption{式で計算した位置は、やってみようの表とぴったり合う。}
\end{bkcenter}

式があれば、表にない時刻の位置も計算できます。たとえば $t=1.1$ なら、$1.1^2=1.21$ なので $20\times1.21=24.2$ です。
\begin{bkcenter}
\begin{bktable}{colspec={Q[l,m,wd=34mm] *{3}{Q[c,m,wd=22mm]}}, rowsep=1.5pt}
時刻 $t$（s） & 1.1 & 1.01 & 1.001 \\
$t^2$ & 1.21 & 1.0201 & 1.002001 \\
位置 $x=20t^2$（cm） & 24.2 & 20.402 & 20.04002 \\
\end{bktable}
\bktabcaption{表にない時刻の位置を、式で計算した。}
\end{bkcenter}

ここからは、「時刻 $t=1.1$ のときの位置」を $x(1.1)$ と書きます。「エックス・かっこ1.1」と読みます（3-3の $f(3)$ と同じ書き方）。$x(1.1)=24.2$、$x(1.0)=20$ です。かけ算の $x\times1.1$ ではないので、気をつけよう。

これで、1.0秒からのあいだを短くしていけます。たとえば1.0秒から1.1秒までなら、進んだ距離は $\Delta x=x(1.1)-x(1.0)=24.2-20=4.2$、平均の速さは $4.2\div0.1=42$ です。ほかの行も、同じ「あと${}-{}$まえ」と「÷$\Delta t$」で計算します。

\begin{bkcenter}
\begin{bktable}{colspec={Q[c,m,wd=22mm] Q[c,m,wd=14mm] Q[c,m,wd=24mm] Q[c,m,wd=20mm] Q[c,m,wd=26mm]}, rowsep=1.5pt, colsep=3pt}
あいだ & $\Delta t$（s） & あとの位置（cm） & $\Delta x$（cm） & 平均の速さ（cm/s） \\
1.0〜2.0 s & 1 & 80 & 60 & 60 \\
1.0〜1.5 s & 0.5 & 45 & 25 & 50 \\
1.0〜1.1 s & 0.1 & 24.2 & 4.2 & 42 \\
1.0〜1.01 s & 0.01 & 20.402 & 0.402 & 40.2 \\
1.0〜1.001 s & 0.001 & 20.04002 & 0.04002 & 40.02 \\
\end{bktable}
\bktabcaption{1.0秒からのあいだを短くしていくと、平均の速さは40に近づく。}
\end{bkcenter}

平均の速さは、60、50、42、40.2、40.02と、40に近づいていきます。40より大きい分（ずれ）を見ると、20、10、2、0.2、0.02です。$\Delta t$ を $\dfrac{1}{10}$ にすると、ずれも $\dfrac{1}{10}$ になっています。だから、$\Delta t$ を短くすれば、ずれはいくらでも0に近くできます。ずれが消えた先の\qty{40}{cm/s}を、\textbf{1.0秒の瞬間の速さ}と決めます。\bkidx{しゅんかんのはやさ}{瞬間の速さ}

\begin{bkfigure}
\begin{tikzpicture}
\begin{axis}[bkplot, width=82mm, height=50mm, xmin=0, xmax=2.3, ymin=0, ymax=95,
  xtick={0,0.5,1,1.5,2}, ytick={0,20,40,60,80}, xlabel={$t$（s）}, ylabel={$x$（cm）}]
  \addplot[bkcurve, domain=0:2.05] {20*x^2};
  \addplot[bkthin, domain=0.7:2.1] {20+60*(x-1)};
  \addplot[bkthin, domain=0.65:2.15] {20+50*(x-1)};
  \addplot[bkcurve2, domain=0.55:2.2, line width=1.2pt] {20+40*(x-1)};
  \addplot[only marks, mark=*, mark size=1.8pt, bkInk] coordinates {(1,20) (2,80) (1.5,45)};
  \node[anchor=south east, font=\sffamily\small] at (axis cs:1,20) {P};
  % 線の右はしに、少しあけてラベル（3本の高さをずらして重ならないようにする）
  \node[anchor=west, font=\sffamily\scriptsize, text=bkGray] at (axis cs:2.25,88.5) {傾き60};
  \node[anchor=west, font=\sffamily\scriptsize, text=bkGray] at (axis cs:2.25,77.5) {傾き50};
  \node[anchor=west, font=\sffamily\scriptsize\bfseries, text=bkOrange!80!black] at (axis cs:2.25,66.5) {傾き40};
\end{axis}
\end{tikzpicture}
\bkfigcaption{2つめの点を P に近づけると、直線の傾きは 60 → 50 → … → 40 と変わっていく。}
\end{bkfigure}

グラフで見ると、2つめの点を点Pにどんどん近づけることにあたります。直線は、だんだん「点Pでグラフにふれるだけの直線」（オレンジ）に近づき、その傾きは40になります。この直線のことは、5-3でくわしく調べます。

\bksubsection{4}{単位つきの式で、時速と秒速をなおす}

\label{bk:5-1:kansan}速さの単位には、時速（km/h）と秒速（m/s）があります。くらべたり計算したりする前に、単位をそろえます（2-10）。そのとき使うのが、次の「単位つきの数」です。

\begin{bkcenter}
\begin{bktable}{colspec={Q[c,m,wd=26mm] Q[l,m,wd=40mm] Q[l,m,wd=50mm]}, rowsep=1.5pt}
単位つきの数 & 読み方 & なぜその数か \\
\qty{1000}{m/km} & 1 km あたり 1000 m & \qty{1}{km}＝\qty{1000}{m}だから \\
\qty{60}{s/min} & 1分あたり60秒 & 1分は60秒だから \\
\qty{60}{min/h} & 1時間あたり60分 & 1時間は60分だから \\
\qty{3600}{s/h} & 1時間あたり3600秒 & $60\times60=3600$（秒）だから \\
\end{bktable}
\bktabcaption{換算に使う単位つきの数（2-10）}
\end{bkcenter}

\noindent\textbf{時速\qty{72}{km}を、秒速になおす}\par
時速\qty{72}{km}は「1時間あたり\qty{72}{km}」です。これを「1秒あたり何m」にします。2つの段に分けます。

\begin{bkfigure}
\begin{tikzpicture}[x=1mm, y=1mm]
  % 1時間（3600秒）の帯と、72 km（72000 m）の帯。はじめの3こと最後の1こだけかき、あいだは「……」でつなぐ
  \foreach \y/\c/\lab in {13/bkBlue/1秒, 0/bkGreen/\qty{20}{m}} {
    \foreach \x in {22,36,50,108} {\fill[\c!10] (\x,\y) rectangle ++(14,7);
      \draw[\c!70!black, line width=0.7pt] (\x,\y) rectangle ++(14,7);
      \node[font=\sffamily\scriptsize] at (\x+7,\y+3.5) {\lab};}
    \node[font=\sffamily\small, text=bkGray] at (86,\y+3.5) {…………};}
  \fill[bkOrange!35] (22,13) rectangle ++(14,7);\draw[bkOrange!80!black, line width=0.9pt] (22,13) rectangle ++(14,7);
  \node[font=\sffamily\scriptsize] at (29,16.5) {1秒};
  \fill[bkOrange!35] (22,0) rectangle ++(14,7);\draw[bkOrange!80!black, line width=0.9pt] (22,0) rectangle ++(14,7);
  \node[font=\sffamily\scriptsize] at (29,3.5) {\qty{20}{m}};
  \node[anchor=east, font=\sffamily\footnotesize] at (20,16.5) {時間};
  \node[anchor=east, font=\sffamily\footnotesize] at (20,3.5) {距離};
  \node[anchor=south, font=\sffamily\footnotesize] at (72,21) {1時間＝3600秒（1秒が3600こ）};
  \node[anchor=north, font=\sffamily\footnotesize] at (72,-1.5) {\qty{72}{km}＝\qty{72000}{m}（3600こに等しく分ける → 1こ $72000\div3600=\qty{20}{m}$）};
\end{tikzpicture}
\bkfigcaption{1時間を3600この1秒に分けると、\qty{72000}{m}も3600こに分かれる。1秒あたり\qty{20}{m}。}
\end{bkfigure}

\textbf{①}\enspace km を m にする。「1 km あたり 1000 m」をかけます。分子の km と分母の km が約分されて消えます。
\[ 72\,\frac{\cancel{\mathrm{km}}}{\mathrm{h}}\times1000\,\frac{\mathrm{m}}{\cancel{\mathrm{km}}}=72000\,\frac{\mathrm{m}}{\mathrm{h}} \]
\textbf{②}\enspace「1時間あたり」を「1秒あたり」にする。1時間は3600秒なので、\qty{3600}{s/h}でわります。わり算は、ひっくり返した数（逆数）をかけるのと同じなので（2-4）、$\div3600\,\mathrm{s/h}$ は $\times\dfrac{1}{3600}\,\mathrm{h/s}$ です。h が約分されて消えます。
\[ 72000\,\frac{\mathrm{m}}{\cancel{\mathrm{h}}}\times\frac{1}{3600}\,\frac{\cancel{\mathrm{h}}}{\mathrm{s}}=20\,\frac{\mathrm{m}}{\mathrm{s}} \]
まとめて1行で書くと、次のようになります。
\[ 72\,\mathrm{km/h}\times1000\,\mathrm{m/km}\div3600\,\mathrm{s/h}=20\,\mathrm{m/s} \]

\noindent\textbf{なぜ「÷3.6」でよいのか}\par
上の式の数だけを見ると、$72\times1000\div3600$ です。
\[ 1000\div3600=\frac{1000}{3600}=\frac{1}{3.6} \]
なので、$72\times1000\div3600=72\div3.6$ です。だから、\textbf{時速を秒速にするには÷3.6}、反対に\textbf{秒速を時速にするには×3.6}になります（2-10）。÷3.6 は近道ですが、どちらにかけるか迷ったら、単位つきの式にもどれば必ずわかります。

\noindent\textbf{秒速を時速にする（反対向き）}\par
秒速\qty{15}{m}を時速になおすときは、①s を h にするために\qty{3600}{s/h}をかけ（s が約分）、②m を km にするために\qty{1000}{m/km}でわります（m が約分）。
\[ 15\,\mathrm{m/s}\times3600\,\mathrm{s/h}\div1000\,\mathrm{m/km}=54\,\mathrm{km/h} \]
$15\times3.6=54$ と同じ答えです。

\bksubsection{5}{スピードメーターは「とても短い時間の平均」}

車のスピードメーターは、タイヤ（またはタイヤにつながった軸）が回る回数を数えています。タイヤが1回転すると、車はタイヤのまわりの長さだけ進みます。

たとえば、タイヤのまわりの長さが\qty{2}{m}（1回転あたり\qty{2}{m}）で、0.5秒のあいだに7回転したとします。進んだ距離は $\qty{2}{m}\times7=\qty{14}{m}$ なので、
\[ \text{平均の速さ}=\frac{14\,\mathrm{m}}{0.5\,\mathrm{s}}=28\,\mathrm{m/s} \]
時速になおすと、
\[ 28\,\mathrm{m/s}\times3600\,\mathrm{s/h}\div1000\,\mathrm{m/km}=100.8\,\mathrm{km/h} \]

\begin{bkfigure}
\begin{tikzpicture}[x=1mm, y=1mm]
  % 第7回：右へころがるタイヤは時計回り。左は、1回転で「まわりの長さ」だけ進むようすを、
  %   実際の長さの比でかく（半径7 mm のタイヤ → 1回転で 2π×7＝43.98 mm 進む）。右は7回転ぶんの帯
  \draw[bkGray, line width=0.8pt] (-10,-7) -- (52,-7);
  \fill[bkInk!20] (0,0) circle (7);
  \draw[bkInk, line width=1pt] (0,0) circle (7);
  \fill[white] (0,0) circle (3.6);
  \draw[bkInk, line width=0.8pt] (0,0) circle (3.6);
  \fill[bkOrange] (0,-5.6) circle (1.1);
  \draw[bk arrow, bkOrange, line width=1pt] (130:10) arc[start angle=130, end angle=50, radius=10];
  \node[font=\sffamily\scriptsize, text=bkOrange!80!black, anchor=south] at (0,10.6) {回る向き};
  \draw[bkInk!45, dash pattern=on 2pt off 1.5pt, line width=0.8pt] (43.98,0) circle (7);
  \fill[bkOrange!55] (43.98,-5.6) circle (1.1);
  \draw[bk arrow, bkBlue] (10,0) -- (33,0);
  \node[font=\sffamily\scriptsize, text=bkBlue, anchor=south] at (21.5,0.8) {進む向き};
  \draw[bkGray, line width=0.6pt, {Bar[width=4pt]}-{Bar[width=4pt]}] (0,-10) -- (43.98,-10);
  \node[font=\sffamily\footnotesize, anchor=north] at (22,-11) {1回転で\qty{2}{m}（タイヤのまわりの長さ）};
  \begin{scope}[shift={(62,0)}]
    \foreach \k in {0,...,6} {\draw[bkGray, fill=bkBlue!8, line width=0.6pt] (\k*9,-7) rectangle ++(9,4);
      \node[font=\sffamily\scriptsize] at (\k*9+4.5,-5) {2 m};}
    \draw[bk arrow, bkBlue] (0,1) -- (63,1);
    \node[font=\sffamily\footnotesize, anchor=north] at (31.5,-11) {0.5秒で7回転 → $2\times7=\qty{14}{m}$};
  \end{scope}
\end{tikzpicture}
\end{bkfigure}

0.5秒はとても短いので、そのあいだに速さはほとんど変わりません。だから、この平均の速さを、その瞬間の速さとして表示しているのです。メーターは、この計算を何度もくり返して、表示を新しくしています。スピードメーターが表示する「瞬間の速さ」は、じつは\textbf{とても短い時間の平均の速さ}です。序章（5-0）の「走っていない1時間の距離を、どうやって知ったのか」の答えは、「とても短い時間に進んだ距離から、1時間あたりの量を計算した」です。

\begin{taisetsu}
\textbf{平均の速さ} $=\dfrac{\Delta x}{\Delta t}$\quad（グラフの2点を結ぶ直線の傾き）\\
\textbf{瞬間の速さ}：$\Delta t$ を0にしないで、どんどん短くしたときに、平均の速さが近づく値\\
\textbf{換算}：$\mathrm{km/h}\times1000\,\mathrm{m/km}\div3600\,\mathrm{s/h}=\mathrm{m/s}$（÷3.6 と同じ）
\end{taisetsu}

\begin{kotoba}
\small
\bkword{平均の速さ}{あるあいだに進んだ距離を、かかった時間でわったもの。そのあいだ、ずっと同じ速さで走ったとしたときの速さ。}
\bkword{瞬間の速さ}{あいだをどんどん短くしたときに、平均の速さが近づいていく値。スピードメーターが表示する速さ。}
\bkword{$\Delta t$、$\Delta x$}{時間と位置の変化量。「デルタ・ティー」「デルタ・エックス」と読む。}
\bkword{$x(1.1)$}{時刻 $t=1.1$ のときの位置 $x$。かけ算ではない。}
\end{kotoba}

\begin{bkfigure}
\begin{tikzpicture}[x=1mm, y=1mm]
  % この章の流れ。箱どうしを矢印でつなぐ
  \node[bk node, text width=30mm, align=center, font=\sffamily\small] (a) at (0,0) {あるあいだの\\平均の速さ $\dfrac{\Delta x}{\Delta t}$};
  \node[bk node, text width=30mm, align=center, font=\sffamily\small] (b) at (42,0) {2点を結ぶ\\直線の傾き};
  \node[bk node, text width=30mm, align=center, font=\sffamily\small] (c) at (84,0) {あいだを\\どんどん短く};
  \node[bk node, text width=30mm, align=center, font=\sffamily\small, fill=bkOrange!12] (d) at (84,-24) {瞬間の速さ\\（近づく先）};
  \node[bk node, text width=34mm, align=center, font=\sffamily\small] (e) at (36,-24) {スピードメーター\\（とても短い時間）};
  \draw[bk link] (a) -- (b);
  \draw[bk link] (b) -- (c);
  \draw[bk link] (c) -- (d);
  \draw[bk link] (d) -- (e);
\end{tikzpicture}
\bkfigcaption{この章の流れ。平均の速さから、瞬間の速さへ。}
\end{bkfigure}

%% ---------------- 例題 ----------------
\bksection{reidai}{例題でたしかめよう}

\begin{reidai}[表から平均の速さを求める]
\bkhatenaref の車の位置を、0.5時間ごとに記録しました。1時間から1.5時間までの平均の速さを求めましょう。
\begin{bkcenter}
\begin{bktable}{colspec={Q[l,m,wd=24mm] *{5}{Q[c,m,wd=11mm]}}, rowsep=1pt}
$t$（h） & 0 & 0.5 & 1.0 & 1.5 & 2.0 \\
$x$（km） & 0 & 30 & 45 & 70 & 100 \\
\end{bktable}
\end{bkcenter}
\kangaekata 時間と位置の変化量を「あと${}-{}$まえ」で求めて、わり算します。\par
\shiki $\Delta t=1.5-1.0=0.5$\quad（h）\par
\hspace*{2.6\zw}$\Delta x=70-45=25$\quad（km）\par
\hspace*{2.6\zw}$\dfrac{\Delta x}{\Delta t}=\dfrac{25}{0.5}=50$\par
\kotae \qty{50}{km/h}
\end{reidai}

\begin{jibun}
例題1の表で、次のあいだの平均の速さを求めましょう。\quad（1）\mbox{1.5〜\qty{2.0}{h}}\quad（2）\mbox{0〜\qty{1.0}{h}}\quad（3）\mbox{0.5〜\qty{1.5}{h}}
\bkjans{（1）\nobreak \qty{60}{km/h}\quad（2）\nobreak \qty{45}{km/h}\quad（3）\nobreak \qty{40}{km/h}}
\end{jibun}

\begin{reidai}[グラフから平均の速さを求める]
\noindent\begin{minipage}[c]{0.55\linewidth}
右の $x$-$t$ グラフの点A（\qty{2}{s}、\qty{6}{m}）から点B（\qty{5}{s}、\qty{21}{m}）までの、平均の速さを求めましょう。\par\vspace{2pt}
\kangaekata 平均の速さは、AとBを結ぶ直線の傾きです。たての変化量（位置）を、横の変化量（時間）でわります。\par
\shiki $\Delta t=5-2=3$\quad（s）、\quad$\Delta x=21-6=15$\quad（m）\par
\hspace*{2.6\zw}$\dfrac{\Delta x}{\Delta t}=\dfrac{15}{3}=5$\par
\kotae \qty{5}{m/s}
\end{minipage}\hfill
\begin{minipage}[c]{0.42\linewidth}
\begin{tikzpicture}
\begin{axis}[bkplot, width=49mm, height=40mm, xmin=0, xmax=6.4, ymin=0, ymax=27, xtick={0,1,...,6}, ytick={0,6,12,18,24},
  xlabel={$t$（s）}, ylabel={$x$（m）}]
  \addplot[bkcurve, domain=0:6, line width=1.3pt] {x^2/2+x+2};
  \addplot[only marks, mark=*, mark size=1.8pt, bkInk] coordinates {(2,6) (5,21)};
  \addplot[bkcurve2, line width=1pt, domain=1.4:5.6] {6+5*(x-2)};
  \node[anchor=north west, font=\sffamily\scriptsize] at (axis cs:2,6) {A};
  \node[anchor=south east, font=\sffamily\scriptsize] at (axis cs:5,21) {B};
\end{axis}
\end{tikzpicture}
\end{minipage}
\end{reidai}

\begin{jibun}
\noindent\begin{minipage}[c]{0.55\linewidth}
右の $x$-$t$ グラフの上の3点、C（\qty{1}{s}、\qty{4}{m}）、D（\qty{4}{s}、\qty{16}{m}）、E（\qty{6}{s}、\qty{26}{m}）について、2点を結ぶ直線の傾きを求め、次の平均の速さを答えましょう。\par
（1）CからD\quad（2）DからE
\end{minipage}\hfill
\begin{minipage}[c]{0.42\linewidth}
\begin{tikzpicture}
\begin{axis}[bkplot, width=49mm, height=40mm, xmin=0, xmax=6.8, ymin=0, ymax=30, xtick={0,1,...,6}, ytick={0,8,16,24},
  xlabel={$t$（s）}, ylabel={$x$（m）}]
  \addplot[bkcurve, domain=0:6.5, line width=1.3pt] {0.2*x^2+3*x+0.8};
  \addplot[only marks, mark=*, mark size=1.8pt, bkInk] coordinates {(1,4) (4,16) (6,26)};
  \node[anchor=north west, font=\sffamily\scriptsize] at (axis cs:1,4) {C};
  \node[anchor=north west, font=\sffamily\scriptsize] at (axis cs:4,16) {D};
  \node[anchor=north west, font=\sffamily\scriptsize] at (axis cs:6,26) {E};
\end{axis}
\end{tikzpicture}
\end{minipage}
\bkjans{（1）\nobreak \qty{4}{m/s}（$\dfrac{16-4}{4-1}=4$）\quad（2）\nobreak \qty{5}{m/s}（$\dfrac{26-16}{6-4}=5$）}
\end{jibun}

\begin{reidai}[あいだを短くして、瞬間の速さを見積もる]
$x=20t^2$（cm）で動くボールの、2.0秒の瞬間の速さを見積もりましょう。\par
\kangaekata 2.0秒から、$\Delta t=0.1$ と $0.01$ のあいだの平均の速さを計算し、近づいていく値を読みます。はじめに、2.0秒の位置を求めます。\par
\shiki $x(2.0)=20\times2.0^2=20\times4=80$\par
\hspace*{2.6\zw}$\Delta t=0.1$ のとき：$x(2.1)=20\times2.1^2=20\times4.41=88.2$\par
\hspace*{2.6\zw}\phantom{$\Delta t=0.1$ のとき：}$\dfrac{88.2-80}{0.1}=\dfrac{8.2}{0.1}=82$\par
\hspace*{2.6\zw}$\Delta t=0.01$ のとき：$x(2.01)=20\times2.01^2=20\times4.0401=80.802$\par
\hspace*{2.6\zw}\phantom{$\Delta t=0.01$ のとき：}$\dfrac{80.802-80}{0.01}=\dfrac{0.802}{0.01}=80.2$\par
\begin{bkcenter}
\begin{bktable}{colspec={Q[c,m,wd=26mm] *{3}{Q[c,m,wd=18mm]}}, rowsep=1pt}
$\Delta t$ & 0.1 & 0.01 & → \\
平均の速さ & 82 & 80.2 & 80 \\
80からのずれ & 2 & 0.2 & 0 \\
\end{bktable}
\end{bkcenter}
\kotae 82、80.2と80に近づくので、およそ\qty{80}{cm/s}
\end{reidai}

\begin{jibun}
例題3のボールについて、$\Delta t=0.1$ と $0.01$ のあいだの平均の速さを計算して、次の時刻の瞬間の速さを見積もりましょう。\par
（1）0.5秒\quad（2）1.5秒
\bkjans{（1）\nobreak $x(0.5)=5$。$\dfrac{7.2-5}{0.1}=22$、$\dfrac{5.202-5}{0.01}=20.2$。およそ\qty{20}{cm/s}\quad（2）\nobreak $x(1.5)=45$。$\dfrac{51.2-45}{0.1}=62$、$\dfrac{45.602-45}{0.01}=60.2$。およそ\qty{60}{cm/s}}
\end{jibun}

\begin{reidai}[時速と秒速をなおす]
（1）時速\qty{72}{km}は、秒速\mbox{何 m}？\quad（2）秒速\qty{15}{m}は、時速\mbox{何 km}？\par
\kangaekata 単位つきの数 \qty{1000}{m/km} と \qty{3600}{s/h} をかけたりわったりして、消したい単位を約分します。\par
\shiki（1）km を消すために \qty{1000}{m/km} をかけ、h を消すために \qty{3600}{s/h} でわる。\par
\hspace*{4.3\zw}$72\,\mathrm{km/h}\times1000\,\mathrm{m/km}\div3600\,\mathrm{s/h}$\par
\hspace*{4.3\zw}$=72\times1000\div3600\ \mathrm{m/s}=20\,\mathrm{m/s}$\par
\hspace*{2.6\zw}（2）s を消すために \qty{3600}{s/h} をかけ、m を消すために \qty{1000}{m/km} でわる。\par
\hspace*{4.3\zw}$15\,\mathrm{m/s}\times3600\,\mathrm{s/h}\div1000\,\mathrm{m/km}$\par
\hspace*{4.3\zw}$=15\times3600\div1000\ \mathrm{km/h}=54\,\mathrm{km/h}$\par
\kotae （1）\qty{20}{m/s}\quad（2）\qty{54}{km/h}
\end{reidai}

\begin{jibun}
（1）時速\qty{90}{km}は秒速\mbox{何 m}でしょう。\quad（2）秒速\qty{12}{m}は時速\mbox{何 km}でしょう。単位つきの式で計算し、÷3.6・×3.6 でも確かめよう。
\bkjans{（1）\nobreak \qty{25}{m/s}\quad（2）\nobreak \qty{43.2}{km/h}}
\end{jibun}

\begin{reidai}[タイヤの回転から速さを求める]
まわりの長さ\qty{1.5}{m}のタイヤが、0.25秒に4回転しました。秒速\mbox{何 m}、時速\mbox{何 km}でしょう。\par
\kangaekata 1回転で、タイヤのまわりの長さ\qty{1.5}{m}だけ進みます。進んだ距離を求めてから、時間でわります。\par
\shiki 進んだ距離：$\qty{1.5}{m}\times4=\qty{6}{m}$\par
\hspace*{2.6\zw}秒速：$\dfrac{6\,\mathrm{m}}{0.25\,\mathrm{s}}=24\,\mathrm{m/s}$\par
\hspace*{2.6\zw}時速：$24\,\mathrm{m/s}\times3600\,\mathrm{s/h}\div1000\,\mathrm{m/km}=86.4\,\mathrm{km/h}$\par
\begin{bkcenter}
\begin{tikzpicture}[x=1mm, y=1mm]
  % 1回転ごとに 1.5 m。4回転で 6 m（長さ 1 m＝16 mm）
  \foreach \k in {0,...,3} {\fill[bkBlue!12] (\k*24,0) rectangle ++(24,5); \draw[bkBlue!70!black, line width=0.6pt] (\k*24,0) rectangle ++(24,5);
    \node[font=\sffamily\scriptsize] at (\k*24+12,2.5) {\the\numexpr\k+1\relax 回転め \qty{1.5}{m}};}
  \draw[bkGray, line width=0.6pt, {Bar[width=4pt]}-{Bar[width=4pt]}] (0,-3) -- (96,-3);
  \node[anchor=north, font=\sffamily\footnotesize] at (48,-4) {0.25秒で $\qty{1.5}{m}\times4=\qty{6}{m}$};
\end{tikzpicture}
\end{bkcenter}
\kotae \qty{24}{m/s}、\qty{86.4}{km/h}
\end{reidai}

\begin{jibun}
まわりの長さ\qty{2.0}{m}のタイヤが、次のように回転しました。それぞれ、秒速\mbox{何 m}、時速\mbox{何 km}でしょう。\par
（1）0.4秒に5回転\quad（2）1秒に10回転
\bkjans{（1）\nobreak $\qty{2.0}{m}\times5=\qty{10}{m}$、$10\div0.4=\qty{25}{m/s}$、\qty{90}{km/h}\quad（2）\nobreak \qty{20}{m/s}、\qty{72}{km/h}}
\end{jibun}

\begin{reidai}[分速を、時速と秒速になおす]
電車が2分で\qty{2.4}{km}進みました。（1）1分あたり\mbox{何 km}（km/min）でしょう。（2）時速\mbox{何 km}でしょう。（3）秒速\mbox{何 m}でしょう。\par
\kangaekata 1時間は60分なので \qty{60}{min/h}、1分は60秒なので \qty{60}{s/min} を使います。\par
\shiki（1）$\dfrac{2.4\,\mathrm{km}}{2\,\mathrm{min}}=1.2\,\mathrm{km/min}$\par
\hspace*{2.6\zw}（2）min を消すために \qty{60}{min/h} をかける。\par
\hspace*{4.3\zw}$1.2\,\mathrm{km/min}\times60\,\mathrm{min/h}=72\,\mathrm{km/h}$\par
\hspace*{2.6\zw}（3）km を m に、min を s にする。\par
\hspace*{4.3\zw}$1.2\,\mathrm{km/min}\times1000\,\mathrm{m/km}\div60\,\mathrm{s/min}=20\,\mathrm{m/s}$\par
\begin{bkcenter}
\begin{tikzpicture}[x=1mm, y=1mm]
  % 時間の帯（60分＝100 mm、目もり1つが1分）。はじめの2分に色
  \fill[bkOrange!50] (0,6) rectangle ({2*100/60},11);
  \foreach \k in {1,...,59} {\draw[bkGray!60, line width=0.25pt] ({\k*100/60},6) -- ({\k*100/60},11);}
  \draw[bkInk, line width=0.6pt] (0,6) rectangle (100,11);
  \node[anchor=south west, font=\sffamily\scriptsize] at (0,11.8) {はじめの2分（色の部分）で\qty{2.4}{km}、1分あたり\qty{1.2}{km}};
  \draw[bkGray, line width=0.6pt, {Bar[width=4pt]}-{Bar[width=4pt]}] (0,3) -- (100,3);
  \node[anchor=north, font=\sffamily\scriptsize] at (50,1.8) {60分（1時間）では $\qty{1.2}{km/min}\times\qty{60}{min}=\qty{72}{km}$（目もり1つが1分）};
\end{tikzpicture}
\end{bkcenter}
\kotae （1）\qty{1.2}{km/min}\quad（2）\qty{72}{km/h}\quad（3）\qty{20}{m/s}
\end{reidai}

\begin{jibun}
自転車が4分で\qty{3}{km}進みました。（1）\mbox{何 km/min}でしょう。（2）時速\mbox{何 km}でしょう。（3）秒速\mbox{何 m}でしょう。
\bkjans{（1）\nobreak \qty{0.75}{km/min}\quad（2）\nobreak \qty{45}{km/h}\quad（3）\nobreak \qty{12.5}{m/s}}
\end{jibun}

\begin{reidai}[行きと帰りで速さがちがうときの平均]
片道\qty{30}{km}の道を、行きは時速\qty{30}{km}、帰りは時速\qty{15}{km}で走りました。往復の平均の速さを求めましょう。\par
\kangaekata 平均の速さは「全部の距離÷全部の時間」です。まず、行きと帰りにかかった時間を求めます。\par
\shiki 行きの時間：$30\div30=1$\quad（h）\par
\hspace*{2.6\zw}帰りの時間：$30\div15=2$\quad（h）\par
\hspace*{2.6\zw}全部の距離：$30+30=60$\quad（km）\par
\hspace*{2.6\zw}平均の速さ：$\dfrac{60}{1+2}=\dfrac{60}{3}=20$\par
\begin{bkcenter}
\begin{tikzpicture}[x=1mm, y=1mm]
  \fill[bkGray!20] (10,-1.5) rectangle (100,1.5);
  \node[bk node, font=\sffamily\small] (A) at (0,0) {家};
  \node[bk node, font=\sffamily\small] (B) at (110,0) {町};
  \draw[bk arrow, bkBlue] (14,6) -- (96,6);
  \node[anchor=south, font=\sffamily\scriptsize, text=bkBlue] at (55,6.8) {行き：時速\qty{30}{km}で\qty{30}{km} → 1時間};
  \draw[bk arrow, bkOrange] (96,-6) -- (14,-6);
  \node[anchor=north, font=\sffamily\scriptsize, text=bkOrange!80!black] at (55,-6.8) {帰り：時速\qty{15}{km}で\qty{30}{km} → 2時間};
\end{tikzpicture}
\end{bkcenter}
\kotae \qty{20}{km/h}（$\dfrac{30+15}{2}=22.5$ ではない）
\end{reidai}

\begin{jibun}
片道\qty{12}{km}の道を、行きは時速\qty{12}{km}、帰りは時速\qty{6}{km}で走りました。往復の平均の速さを求めましょう。
\bkjans{行き1時間、帰り2時間。$\dfrac{24}{3}=\qty{8}{km/h}$}
\end{jibun}

\clearpage
%% ---------------- まちがえやすいところ・物理工学 ----------------
\bksection{tsumazuki}{まちがえやすいところをチェック}

\begin{tsumazuki}{平均の速さは「速さの平均」ではない}
\qty{120}{km}はなれた町まで、行きは時速\qty{60}{km}、帰りは時速\qty{40}{km}で走りました。往復の平均の速さは？
\par\vspace{2pt}
\bkng{$\dfrac{60+40}{2}=\qty{50}{km/h}$}\bkok{行き2時間、帰り3時間。$\dfrac{240}{2+3}=\qty{48}{km/h}$}
\par\vspace{2pt}
平均の速さは、いつも「全部の距離÷全部の時間」で求めます。おそい帰りのほうが時間が長いので、平均はおそいほうに引っぱられます。
\end{tsumazuki}

\begin{tsumazuki}{単位をそろえずにわり算する}
$\qty{3}{km}$ を5分で走ったとき、$3\div5=0.6$ は「\qty{0.6}{km/min}」です。「\qty{0.6}{m/s}」ではありません。
\bkng{$3\div5=0.6$ なので \qty{0.6}{m/s}}\bkok{$0.6\,\mathrm{km/min}\times1000\,\mathrm{m/km}\div60\,\mathrm{s/min}=10\,\mathrm{m/s}$}
数だけでなく、単位もいっしょに書いて計算すると、まちがいに気づけます。
\end{tsumazuki}

\begin{tsumazuki}{換算で、かけるとわるを逆にする}
\bkng{$72\,\mathrm{km/h}\times3600\,\mathrm{s/h}=259200$}\bkok{$72\,\mathrm{km/h}\div3600\,\mathrm{s/h}\times1000\,\mathrm{m/km}=20\,\mathrm{m/s}$}
左の式では h が約分されず、$\mathrm{km\cdot s/h^2}$ というおかしな単位が残ります。単位が m/s になるかで確かめよう。
\end{tsumazuki}

\begin{tsumazuki}{$x(1.1)$ をかけ算と読む}
\bkng{$x(1.1)=x\times1.1$}\bkok{$x(1.1)$ は「$t=1.1$ のときの位置」。$x=20t^2$ なら $20\times1.1^2=24.2$}
かっこの前の $x$ は「位置を求めるきまり」の名前で、かけられる数ではありません（3-3の $f(3)$ と同じ）。
\end{tsumazuki}

\begin{tsumazuki}{瞬間の速さを $0\div0$ で求めようとする}
\bkng{あいだを0にして、$\dfrac{0}{0}$ を計算する}\bkok{あいだを0.1、0.01、0.001と短くして、近づく先を読む}
0でわる計算はできません（1-14）。あいだは0にしないで、短くしていきます。
\end{tsumazuki}

\begin{butsuri}
\begin{itemize}[leftmargin=1.2\zw, itemsep=2pt, topsep=0pt, label={\color{bkPurple}●}]
\item \textbf{区間の平均速度の取りしまり}\par
外国の高速道路などには、2つの地点を通った時刻から、そのあいだの平均の速さを計算して、速すぎる車を見つけるしくみがあります。
\item \textbf{スマートフォンの速さ}\par
地図アプリは、短い時間ごとの位置（GPS）の変化などから、速さを計算して表示しています。
\item \textbf{100 m 走の記録}\par
9.58秒の世界記録の平均の速さは、$100\div9.58$ で約\qty{10.4}{m/s}。でも、速さがいちばん大きいところでは約\qty{12}{m/s}あったと計算されています。スタートの直後はおそいので、平均はそれより小さくなるのです。
\end{itemize}
\end{butsuri}

\begin{motto}[2]
\textbf{もどるときの速さはマイナス}\par
ボールが坂をのぼって、また下りてくるとき、下りるあいだは位置 $x$ が小さくなるので、$\Delta x$ はマイナスになります（2-2：マイナスは反対の向き）。たとえば2秒で\qty{6}{m}もどったら、$\dfrac{\Delta x}{\Delta t}=\dfrac{-6}{2}=\qty{-3}{m/s}$。このように向きもふくめた速さを\textbf{速度}といいます。$x$-$t$ グラフでは、右下がりの直線の傾きにあたります。速度は5-5でくわしく使います。\bkidx{そくど}{速度}
\end{motto}

\clearpage
%% ---------------- 章末問題 ----------------
\bkendproblems
\bkproblemgroup{基本}{2}

\bkproblem 電車の位置を2分ごとに記録しました。
\begin{bkcenter}
\begin{bktable}{colspec={Q[l,m,wd=24mm] *{4}{Q[c,m,wd=12mm]}}, rowsep=1pt}
$t$（min） & 0 & 2 & 4 & 6 \\
$x$（km） & 0 & 1.2 & 4.0 & 6.0 \\
\end{bktable}
\end{bkcenter}
（1）0〜6分の平均の速さは、\mbox{何 km/min} でしょう。\quad（2）2〜4分の平均の速さは、\mbox{何 km/min} でしょう。\par
（3）（1）の速さは、時速\mbox{何 km}でしょう。単位つきの式で求めましょう。
\bkpans{（1）\nobreak \qty{1.0}{km/min}\quad（2）\nobreak \qty{1.4}{km/min}\quad（3）\nobreak $1.0\,\mathrm{km/min}\times60\,\mathrm{min/h}=\qty{60}{km/h}$}

\bkproblem $x$-$t$ グラフの上に、3つの点A（\qty{0}{s}、\qty{0}{m}）、B（\qty{4}{s}、\qty{12}{m}）、C（\qty{10}{s}、\qty{24}{m}）があります。AからB、BからC、AからCの平均の速さを、それぞれ求めましょう。
\bkpans{A→B \qty{3}{m/s}、B→C \qty{2}{m/s}、A→C \qty{2.4}{m/s}}

\bkproblem 単位つきの式で求めましょう。（1）時速\qty{36}{km}は秒速\mbox{何 m}？\quad（2）秒速\qty{5}{m}は時速\mbox{何 km}？\quad（3）時速\qty{300}{km}の新幹線は、秒速\mbox{何 m}？（小数第1位まで）
\bkpans{（1）\nobreak \qty{10}{m/s}\quad（2）\nobreak \qty{18}{km/h}\quad（3）\nobreak \qty{83.3}{m/s}}

\bkproblem タイヤのまわりの長さが\qty{1.8}{m}の車で、タイヤが0.2秒に3回転しました。このときの速さは、秒速\mbox{何 m}ですか。時速\mbox{何 km}ですか。
\bkpans{\qty{27}{m/s}、\qty{97.2}{km/h}}

\bkproblem $x=20t^2$（cm）で動くボールの、2.5秒の瞬間の速さを見積もりましょう。$\Delta t=0.1$ と $0.01$ のあいだの平均の速さを計算して、近づく値を答えましょう。
\bkpans{$x(2.5)=125$、$x(2.6)=135.2$ より 102、$x(2.51)=126.002$ より 100.2。およそ\qty{100}{cm/s}}

\bkproblemgroup{チャレンジ}{2}

\bkproblem 位置が $x=5t^2$（m）で表される物体があります。
（1）3.0秒から3.1秒までの平均の速さを求めましょう。\quad（2）3.0秒から3.01秒までの平均の速さを求めましょう。\par
（3）3.0秒の瞬間の速さは、およそいくつと考えられますか。
\bkpans{（1）\nobreak \qty{30.5}{m/s}（$\dfrac{48.05-45}{0.1}$）\quad（2）\nobreak \qty{30.05}{m/s}\quad（3）\nobreak およそ\qty{30}{m/s}}

\bkproblem 片道\qty{6}{km}の山道を、上りは時速\qty{4}{km}、下りは時速\qty{6}{km}で歩きました。往復の平均の速さを求めましょう。
\bkpans{上り1.5時間、下り1時間。$\dfrac{12}{2.5}=\qty{4.8}{km/h}$}

\bkproblem 秒速\qty{340}{m}で進む音は、時速\mbox{何 km}でしょう。単位つきの式で求めましょう。また、時速\qty{1000}{km}で飛ぶ飛行機と、どちらが速いでしょう。
\bkpans{$340\,\mathrm{m/s}\times3600\,\mathrm{s/h}\div1000\,\mathrm{m/km}=\qty{1224}{km/h}$。音のほうが速い}

\begin{furikaeri}
\begin{checklist}
\item 平均の速さを $\dfrac{\Delta x}{\Delta t}$ で求められる。\hfill{\footnotesize\sffamily\color{bkTheme}→\pageref{bk:5-1:heikin}ページ}
\item 平均の速さが、2点を結ぶ直線の傾きだと説明できる。\hfill{\footnotesize\sffamily\color{bkTheme}→\pageref{bk:5-1:katamuki}ページ}
\item あいだを短くして、瞬間の速さを見積もれる。\hfill{\footnotesize\sffamily\color{bkTheme}→\pageref{bk:5-1:shunkan}ページ}
\end{checklist}
\end{furikaeri}

\bkfukidashi{mina}{瞬間の速さは「近づいていく値」なんだね。次の章では、近づいていくようすを式で追いかけてみよう。}
